% Notas de aula — Capítulo 7: Processos de Precipitação
% Curso: Meteorologia Prática (baseado em R. Stull, Practical Meteorology, CC BY-NC-SA 4.0)
% Caixas verdes = conteúdo literal do quadro; linhas verdes = encenação.
\documentclass[11pt]{article}
\usepackage{../comum/preambulo}
\graphicspath{{../figuras/cap07/}}

\title{Capítulo 7 --- Processos de Precipitação\\
{\large Notas de aula (roteiro de quadro-negro)}}
\author{Curso de Meteorologia Prática}
\date{}

\begin{document}
\maketitle
\thispagestyle{fancy}

\noindent\textbf{Plano sugerido:} 2 aulas de 2\,h.
\emph{Aula 1:} \S\ref{sec:problema}--\S\ref{sec:kohler} (o problema das
escalas, Kelvin e Köhler, nucleação).
\emph{Aula 2:} \S\ref{sec:crescimento}--\S\ref{sec:distribuicao}
(crescimento por difusão e por coleta, Bergeron, distribuição de
tamanhos).

\section{O problema central: 6 ordens de grandeza}\label{sec:problema}

\begin{noquadro}[a escada que organiza a aula]
\[ \underbrace{0{,}1\,\mu\text{m}}_{\text{CCN}} \to
   \underbrace{10\,\mu\text{m}}_{\text{gotícula de nuvem}} \to
   \underbrace{1\,\text{mm}}_{\text{gota de chuva}} \]
gotícula $\to$ gota: volume cresce $\mathbf{10^6}\times$\quad
— e a nuvem faz isso em $\sim$30\,min.
\end{noquadro}

O Cap.~6 terminou com o ar saturado e condensando
\javimos{6}{caminhos para a saturação}; mas ``condensou'' ainda não é
``choveu''. Explicar como a atmosfera sobe essa escada em meia hora é o
objetivo de hoje. Spoiler para manter a turma orientada: difusão de
vapor constrói a nuvem mas \emph{não} constrói a chuva; a chuva precisa
de colisões (nuvens quentes) ou de gelo (Bergeron). No fim, a
distribuição de gotas resultante é exatamente o que o radar enxerga
\veremos{8}{refletividade de radar}.

\section{Efeito de curvatura: barreira de Kelvin}\label{sec:kelvin}

\quadro{Desenhar gota minúscula com moléculas escapando fácil da
superfície muito curva; ao lado, superfície plana segurando melhor as
moléculas.}

\begin{formadif}[equação de Kelvin]
A tensão superficial $\sigma$ torna mais fácil evaporar de uma
superfície curva. Minimizando a energia livre (superfície vs.\ volume):
\[ \frac{e_s(r)}{e_s(\infty)} =
\exp\!\left(\frac{2\sigma}{\rho_L \Re_v T\, r}\right)
 \equiv \exp\!\left(\frac{A}{r}\right),\qquad
A \simeq 1{,}2\times10^{-3}\ \mu\mathrm{m}\ \ (\text{a }273\ \mathrm{K}). \]
\end{formadif}

\begin{noquadro}[a consequência brutal]
gota de $0{,}01\,\mu$m exige UR $\simeq 112\%$ para sobreviver;\\
a atmosfera raramente passa de $100{,}5\%$\\[2pt]
$\Rightarrow$ \textbf{nuvens não nascem do ar puro}: nascem sobre
núcleos de condensação (CCN)\\
poeira, sal marinho, sulfatos, fumaça — sem aerossol, sem nuvem.
\end{noquadro}

É um dos raros lugares onde a microfísica de superfícies decide o clima
do planeta: a nucleação homogênea é proibida pela exponencial de Kelvin,
então \emph{todo} o céu nublado que existe passou por um grão de
aerossol \veremos{21}{efeitos do aerossol no clima}.

\section{Efeito de soluto e curvas de Köhler}\label{sec:kohler}

\quadro{Construir a curva de Köhler somando os dois efeitos no quadro:
Kelvin ($+A/r$) empurra para cima, Raoult ($-B/r^3$) para baixo; marcar
$r^*$ e $S^*$ no máximo.}

Sais dissolvidos \emph{reduzem} $e_s$ (efeito Raoult — as moléculas de
sal ocupam lugar na superfície), competindo com Kelvin:

\begin{formalivro}[curva de Köhler, Stull \S7.2]
\[ S(r) \simeq 1 + \frac{A}{r} - \frac{B\,r_d^3}{r^3}, \]
com $r_d$ o raio seco do núcleo e $B$ dependente da higroscopicidade
(parametrizada hoje pelo $\kappa$ de Petters--Kreidenweis).
\end{formalivro}

\begin{noquadro}[ativação]
máximo da curva $\Rightarrow$ supersaturação crítica $S^*$ e raio
crítico $r^*$\\
passou de $r^*$: gota \textbf{ativada}, cresce sozinha\\
núcleo maior/mais salgado $\Rightarrow$ $S^*$ menor $\Rightarrow$ ativa
primeiro
\end{noquadro}

Daqui sai o elo aerossol--clima: ar marítimo limpo ativa
$\sim$100\,gotas/cm$^3$; ar continental poluído, $\sim$1000/cm$^3$ —
mais gotas menores, nuvens mais brancas e duradouras (efeito indireto do
aerossol \veremos{21}{forçante de aerossol}) e chuva quente mais lenta
(gotículas pequenas demais para colidir, \S\ref{sec:crescimento}).
Simulação com o pacote \texttt{pyrcel} no Caso~2 do notebook: quem sobe
mais rápido gera $S_{max}$ maior e ativa mais núcleos (exercício N2).

\textbf{Gelo}: nucleação homogênea só abaixo de $-38\,°$C; entre 0 e
$-38\,°$C exige núcleos de gelo (IN) — raríssimos ($1/10^5$ dos CCN),
tipicamente poeiras minerais. Por isso nuvens até $-20\,°$C são
majoritariamente de água \emph{super-resfriada}
\javimos{4}{$e_s$ água × gelo} — fato central para o Bergeron adiante e
para o congelamento de aeronaves.

\section{Crescimento por difusão: rápido no início, lento demais no fim}
\label{sec:crescimento}

\begin{formadif}[difusão de vapor para a gota]
O fluxo difusivo de vapor sobre a gota dá
\[ r\,\frac{dr}{dt} = D\,\frac{\rho_{v,\infty} - \rho_{v,s}}{\rho_L}
\quad\Longrightarrow\quad r(t) = \sqrt{r_0^2 + 2Ct},\qquad
r \propto \sqrt{t}. \]
A raiz quadrada é a assinatura de todo crescimento limitado por difusão
(mesma matemática da camada de difusão em eletroquímica).
\end{formadif}

\begin{exemplo}[o gargalo da chuva quente]
Com supersaturação típica de nuvem ($S - 1 \simeq 0{,}5$\,\%), a
gotícula vai de 1 a 10\,$\mu$m em minutos — mas de 10\,$\mu$m a 1\,mm
levaria \emph{dias} (o $\sqrt{t}$ mata: cada década de raio custa 100
vezes mais tempo). A difusão constrói a nuvem, mas \textbf{não} constrói
a chuva. \emph{Verificação no Caso~3 do notebook.}
\end{exemplo}

\subsection{Colisão--coalescência (chuva quente)}

\quadro{Desenhar a gota grande caindo e varrendo as pequenas no
caminho — o ``coletor'' cresce como bola de neve.}

Gotas maiores caem mais rápido e varrem as menores:

\begin{formalivro}[crescimento por coleta]
\[ \frac{dr}{dt} = \frac{E\,\bar{L}\,\Delta w_T}{4\rho_L}, \]
com $E$ eficiência de coleta, $\bar L$ conteúdo de água líquida e
$\Delta w_T$ a diferença de velocidades terminais. Crescimento
\emph{acelera} com $r$ — é a avalanche que faltava: de 20\,$\mu$m a
1\,mm em $\sim$20--30\,min em nuvens tropicais.
\end{formalivro}

Note a inversão de regime: difusão desacelera, coleta acelera; o
cruzamento (por volta de 20\,$\mu$m) é o gargalo que decide se a nuvem
chove — e é onde o número de gotas (Köhler!) entra: nuvem poluída, com
muitas gotas pequenas, demora a atravessá-lo. O exercício N3 acha esse
raio de cruzamento.

\subsection{Processo de Bergeron (nuvens frias)}

\quadro{Retomar a figura das duas curvas $e_s$ do Cap.~4; marcar o
máximo da diferença em $-12\,°$C.}

Num mar de gotículas super-resfriadas, o raro cristal de gelo vê o
ambiente \emph{supersaturado} (porque $e_{s,\text{gelo}} <
e_{s,\text{água}}$ \javimos{4}{as duas curvas}): cresce por difusão às
custas das gotículas — que evaporam para alimentá-lo —, depois agrega
(flocos) e frag\-menta (multiplicação de gelo). É o mecanismo dominante
da precipitação de latitudes médias: a maior parte da ``chuva'' nasce
neve e derrete ao cair. A fase em que chega ao chão depende do perfil de
$T_w$ \javimos{4}{bulbo úmido} — chuva, neve, ou a temida chuva
congelante \veremos{14}{granizo, o primo convectivo}.

\section{Velocidade terminal e distribuição de tamanhos}
\label{sec:distribuicao}

\begin{formalivro}[velocidades terminais]
Gotículas (Stokes, $r < 40\,\mu$m): $w_T = k_1 r^2$,
$k_1 = 1{,}19\times10^{8}$\,m$^{-1}$s$^{-1}$;
gotas de chuva ($\sim$0{,}6--2\,mm): $w_T \simeq$ 4--9\,m/s (regime
turbulento; o T4 deduz Stokes igualando peso e arrasto viscoso).
\end{formalivro}

\begin{formalivro}[Marshall--Palmer]
\[ N(D) = N_0\,e^{-\Lambda D},\qquad N_0 = 8\times10^{6}\ \mathrm{m^{-4}},
\quad \Lambda = 4{,}1\,R^{-0{,}21}\ \mathrm{mm^{-1}}\ (R\ \text{em mm/h}). \]
Chuva mais intensa $\Rightarrow$ $\Lambda$ menor $\Rightarrow$ gotas
grandes mais prováveis.
\end{formalivro}

\begin{noquadro}[a ponte para o radar]
da Marshall--Palmer saem:\\
taxa de chuva $R \propto \int N\,w_T\,D^3\,dD$ \qquad refletividade
$Z \propto \int N\,D^6\,dD$\\[2pt]
o $D^6$ favorece brutalmente as gotas grandes $\Rightarrow$ relação
$Z$--$R$ do Cap.~8 (N4).
\end{noquadro}

\section{Medindo a precipitação}

Fecho instrumental (Stull \S7.8): \textbf{pluviômetro} (o padrão-ouro
pontual; o de báscula dá a série temporal), \textbf{disdrômetro}
(mede a distribuição $N(D)$ — a Marshall--Palmer observada
diretamente!), medida de neve (relação $\sim$10\,mm de neve por 1\,mm
de água, muito variável), e as estimativas de área por radar
($Z$--$R$ \veremos{8}{radar}) e satélite — que precisam dos
pluviômetros para calibrar. A intensidade classifica-se como fraca
($<2{,}5$), moderada (2,5--7,6) e forte ($>7{,}6$\,mm/h).

\section{Síntese da aula}

\begin{noquadro}[mapa final --- canto do quadro]
\[ \text{CCN} \xrightarrow{\text{Köhler}} \text{gotícula}
   \xrightarrow{\text{difusão } \sqrt{t}} 10\,\mu\text{m}
   \xrightarrow[\text{ou Bergeron}]{\text{colisões}} \text{gota}
   \xrightarrow{\text{Marshall--Palmer}} \text{radar (Cap.~8)} \]
sem aerossol não há nuvem; sem colisões ou gelo não há chuva.
\end{noquadro}

\section{Exercícios complementares}\label{sec:exercicios}

Soluções (professor) em \texttt{cap07\_solucoes.ipynb}.

\subsection*{Teóricos}

\begin{enumerate}[label=\textbf{T\arabic*.},itemsep=4pt]
  \item Derive a equação de Kelvin minimizando a energia livre
        $G = -n\,\Delta\mu\,\tfrac43\pi r^3/v_m + 4\pi r^2\sigma$ e
        calcule a UR de equilíbrio para gotas de 0{,}01, 0{,}1 e
        1\,$\mu$m.
  \item Para a curva de Köhler $S = 1 + A/r - Br_d^3/r^3$, mostre que
        $r^* = \sqrt{3Br_d^3/A}$ e $S^* = 1 + \sqrt{4A^3/(27Br_d^3)}$;
        interprete a dependência com $r_d$.
  \item Integre $r\,dr/dt = C$ e mostre que o tempo para ir de $r_1$ a
        $r_2$ por difusão é $t = (r_2^2 - r_1^2)/2C$; com o $C$ típico
        de nuvem, estime os tempos 1$\to$10\,$\mu$m e
        10\,$\mu$m$\to$1\,mm.
  \item Derive a velocidade de Stokes
        $w_T = \tfrac{2}{9}\,g r^2(\rho_L-\rho)/\mu$ igualando peso e
        arrasto viscoso, e verifique o $k_1$ do livro.
  \item Da Marshall--Palmer, mostre que o conteúdo de água da chuva é
        $M = \pi\rho_L N_0/(6\Lambda^4)\cdot\Gamma(4)$ e que o diâmetro
        médio ponderado por massa é $D_m = 4/\Lambda$.
\end{enumerate}

\subsection*{Numéricos (no notebook)}

\begin{enumerate}[label=\textbf{N\arabic*.},itemsep=4pt]
  \item Trace famílias de curvas de Köhler variando $\kappa$ (0{,}1 a
        1{,}3) e $r_d$; ache numericamente $S^*(r_d)$ e compare com a
        fórmula do T2.
  \item Com o \texttt{pyrcel} (ou o modelo simplificado do notebook),
        varra a velocidade de subida $W$ de 0{,}1 a 5\,m/s e trace a
        fração de aerossol ativada e o pico de supersaturação
        $S_{max}(W)$.
  \item Integre o crescimento combinado difusão + coleta e ache o tempo
        total gotícula$\to$gota para $\bar L$ = 0{,}3, 1{,}0 e 3{,}0
        g/m$^3$; identifique o raio em que a coleta assume o comando.
  \item Da Marshall--Palmer, calcule numericamente $R$ (integral de
        $N\,w_T\,D^3$) e $Z$ ($\int N D^6$) para $\Lambda$ de chuvas
        fraca, moderada e forte, e verifique a relação $Z$--$R$ tipo
        $Z = a R^{1.4}$ que o Cap.~8 usará.
\end{enumerate}

\vspace{1em}
\noindent\rule{\linewidth}{0.4pt}\\
{\scriptsize Material derivado de R.~Stull, \emph{Practical Meteorology}
(CC BY-NC-SA 4.0); este documento herda a mesma licença.}

\end{document}
